\chapter{de Rham Cohomology}\label{de-rham}
\chapterstatus{complete}

\section{Introduction}\label{de-rham:section-introduction}

Every exact form is closed, because $d \circ d = 0$. The converse is a question about the shape of the manifold. On $\RR^n$, and on any convex open set, every closed form of
positive degree is exact (the Poincar\'e lemma). On the punctured plane $\RR^2 \setminus \{0\}$ the angle form $(x\,dy - y\,dx)/(x^2 + y^2)$ is closed but not exact: its integral
around the origin is $2\pi$, while the integral of an exact form around a closed curve is $0$. The obstruction comes from the hole at the origin. \emph{De Rham cohomology},
closed forms modulo exact forms, measures such obstructions in every degree:
\[
H^k_{\mathrm{dR}}(M) = \frac{\{\text{closed } k\text{-forms}\}}{\{\text{exact } k\text{-forms}\}} .
\]
The de Rham theorem (Theorem~\ref{de-rham:theorem-de-rham}) says that this analytic invariant is the singular cohomology of $M$ with real coefficients, a topological invariant.
This is the bridge between analysis and topology on which Hodge theory is built: the Hodge theorem selects a canonical harmonic representative in each de Rham class
(Chapter~\ref{hodge-theorem}), and on a complex manifold the type decomposition of harmonic forms gives the Hodge decomposition (Chapter~\ref{kahler}).

The chapter goes as follows.
\begin{enumerate}
\item The de Rham complex, first examples, homotopy invariance with an explicit homotopy operator, and the Poincar\'e lemma with an explicit formula for primitives
(Section~\ref{de-rham:section-complex}).
\item The Mayer--Vietoris sequence, with its connecting homomorphism made explicit, and the computation of the cohomology of spheres, punctured planes and projective spaces.
\item The de Rham theorem, periods of $1$-forms, the K\"unneth formula, cohomology with compact supports (fibre integration, the compact Poincar\'e lemma and Mayer--Vietoris), Poincar\'e duality proved by induction on a good cover,
and the cohomology of compact Lie groups (Section~\ref{de-rham:section-de-rham-theorem}).
\item Degree theory with forms: $\int_M$ identifies $H^n_{\mathrm{dR},c}(M)$ with $\RR$, the degree of a proper map, its computation as a signed count of preimages, homotopy invariance, the fundamental theorem of algebra,
and the comparison with the homological degree (Section~\ref{de-rham:section-degree}).
\end{enumerate}
Currents, the forms with distribution coefficients, which contain both forms and submanifolds and also compute de Rham cohomology, are developed in Chapter~\ref{currents}.
References: \cite{bott-tu}, \cite{warner-foundations}, \cite[Chapter 0]{griffiths-harris}.

\noindent\textbf{Prerequisites.} Chapter~\ref{differential-forms} (Differential Forms and Integration), Section~\ref{sheaves:section-sheaf-cohomology}
(Sheaf Cohomology), and Chapter~\ref{poincare-duality} (Topological Manifolds and Poincar\'e Duality).

\section{The de Rham complex}\label{de-rham:section-complex}

\begin{definition}\label{de-rham:definition-de-rham}
The \emph{de Rham complex} of a smooth manifold $M$ is $0 \to \Omega^0(M) \xrightarrow{d} \Omega^1(M) \xrightarrow{d} \cdots$, a complex by Theorem~\ref{differential-forms:theorem-exterior-derivative}. Its
cohomology $H^k_{\mathrm{dR}}(M) = \ker(d \colon \Omega^k \to \Omega^{k+1})/\operatorname{im}(d \colon \Omega^{k-1} \to \Omega^k)$ is the \emph{de Rham cohomology}; its elements are classes $[\omega]$ of
\emph{closed} forms modulo \emph{exact} ones. The wedge product makes $H^\bullet_{\mathrm{dR}}(M)$ a graded-commutative $\RR$-algebra, and a smooth map $f \colon N \to M$ induces an algebra homomorphism
$f^* \colon H^\bullet_{\mathrm{dR}}(M) \to H^\bullet_{\mathrm{dR}}(N)$. Replacing $\Omega^\bullet(M)$ by the compactly supported forms $\Omega^\bullet_c(M)$ gives $H^\bullet_{\mathrm{dR},c}(M)$, functorial for
open inclusions (extension by zero) and for proper maps.
\end{definition}

\begin{proposition}\label{de-rham:proposition-h0}
$H^0_{\mathrm{dR}}(M) \cong \RR^{\pi_0(M)}$, the locally constant functions. The wedge product is well defined on classes, and $f^*$ commutes with $d$, hence descends to cohomology.
\end{proposition}

\begin{proof}
A function with $df = 0$ has vanishing differential, hence is constant on connected components (Theorem~\ref{real-analysis:theorem-mean-value} in charts, and connectedness). If $\alpha$, $\beta$ are closed then
$d(\alpha \wedge \beta) = 0$, and if $\alpha = d\gamma$ then $\alpha \wedge \beta = d(\gamma \wedge \beta)$, because $d(\gamma \wedge \beta) = d\gamma \wedge \beta \pm \gamma \wedge d\beta$ and $d\beta = 0$; naturality of $d$ is
Theorem~\ref{differential-forms:theorem-exterior-derivative}.
\end{proof}

\begin{example}\label{de-rham:example-first-computations}
\begin{enumerate}
\item \emph{A point.} On the $0$-dimensional manifold $\{p\}$, $\Omega^0 = \RR$ (the functions are the constants) and $\Omega^k = 0$ for $k \geq 1$. So $H^0_{\mathrm{dR}}(\{p\}) = \RR$ and
$H^k_{\mathrm{dR}}(\{p\}) = 0$ for $k \geq 1$; more generally a $0$-dimensional manifold, a countable discrete set $S$, has $H^0 = \RR^S$ and nothing else. By homotopy invariance
(Theorem~\ref{de-rham:theorem-homotopy-invariance} below) every contractible manifold, such as $\RR^n$ or a ball, has the de Rham cohomology of a point: this is the Poincar\'e lemma.
\item \emph{The real line.} Every $1$-form on $\RR$ is $f\,dx$, and it is exact: $f\,dx = dF$ with $F(x) = \int_0^xf$. So $H^0_{\mathrm{dR}}(\RR) = \RR$ and $H^k_{\mathrm{dR}}(\RR) = 0$ for
$k \geq 1$ (there are no forms of degree $\geq 2$).
\item \emph{The circle.} On $S^1 = \RR/2\pi\ZZ$ every $1$-form is closed, and is $f(\theta)\,d\theta$ with $f$ periodic. It is exact if and only if $\int_0^{2\pi}f = 0$: then
$F(\theta) = \int_0^\theta f$ is periodic and $dF = f\,d\theta$; conversely $\int_{S^1}dF = 0$. So $\int_{S^1} \colon H^1_{\mathrm{dR}}(S^1) \to \RR$ is an isomorphism, and $[d\theta]$ is a basis.
Here $d\theta$ is a closed $1$-form which is not the differential of a function, since $\theta$ is not a function on $S^1$.
\item \emph{The punctured plane.} The angle form $\alpha = (x\,dy - y\,dx)/(x^2 + y^2)$ is closed and not exact on $\RR^2 \setminus \{0\}$
(Counterexample~\ref{differential-forms:counterexample-closed-not-exact}); so $H^1_{\mathrm{dR}}(\RR^2 \setminus \{0\}) \neq 0$. It is the pullback of $d\theta$ by
$(x, y) \mapsto (x, y)/r \in S^1$, and we will see that it spans $H^1$ (Example~\ref{de-rham:example-angle-form}).
\end{enumerate}
\end{example}

\begin{theorem}[Homotopy invariance and the Poincar\'e lemma]\label{de-rham:theorem-homotopy-invariance}
Let $i_0, i_1 \colon M \to M \times \RR$ be $p \mapsto (p, 0)$ and $p \mapsto (p, 1)$. There is a linear map $K \colon \Omega^k(M \times \RR) \to \Omega^{k-1}(M)$ with
\[
i_1^*\omega - i_0^*\omega = d(K\omega) + K(d\omega).
\]
Consequently smoothly homotopic maps induce the same map on de Rham cohomology, homotopy equivalent manifolds have isomorphic de Rham cohomology, and $H^k_{\mathrm{dR}}(U) = 0$ for $k > 0$ and
$H^0_{\mathrm{dR}}(U) = \RR$ for every convex open $U \subseteq \RR^n$ (the \emph{Poincar\'e lemma}).
\end{theorem}

\begin{proof}
\emph{The decomposition.} Let $t$ be the coordinate of $\RR$. In local coordinates $x$ on $M$ and $t$, a $k$-form on $M \times \RR$ is a sum of terms $f(x, t)\,dx^I$ and
$g(x, t)\,dt \wedge dx^J$; so it can be written uniquely as
\[
\omega = \alpha_t + dt \wedge \beta_t,
\]
where $\alpha_t$ is a $k$-form and $\beta_t$ a $(k-1)$-form on $M$, depending smoothly on the parameter $t$ (their coefficients are $f(\cdot, t)$ and $g(\cdot, t)$); this is
independent of the coordinates on $M$. For instance $i_s^*\omega = \alpha_s$, since $i_s^*dt = 0$. Write $d_M$ for the exterior derivative in the $M$-directions at fixed $t$, and
$\partial_t$ for the derivative of the coefficients in $t$. Then $d = d_M + dt \wedge \partial_t$ on forms without $dt$, and
\[
d\omega = d_M\alpha_t + dt \wedge \partial_t\alpha_t + d(dt \wedge \beta_t) = d_M\alpha_t + dt \wedge (\partial_t\alpha_t - d_M\beta_t),
\]
because $d(dt \wedge \beta_t) = -dt \wedge d\beta_t = -dt \wedge (d_M\beta_t + dt \wedge \partial_t\beta_t) = -dt \wedge d_M\beta_t$.

\emph{The homotopy operator.} Define $K\omega = \int_0^1\beta_t\,dt$, the $(k-1)$-form on $M$ whose coefficients are the integrals over $t \in [0, 1]$ of those of $\beta_t$. By the
formula for $d\omega$, its $dt$-part is $\partial_t\alpha_t - d_M\beta_t$, so $K(d\omega) = \int_0^1(\partial_t\alpha_t - d_M\beta_t)\,dt$. On the other hand
$d(K\omega) = \int_0^1d_M\beta_t\,dt$, by differentiation under the integral sign (Proposition~\ref{real-analysis:proposition-parameter-integrals}). Adding, and using the fundamental
theorem of calculus for the coefficients,
\[
d(K\omega) + K(d\omega) = \int_0^1\partial_t\alpha_t\,dt = \alpha_1 - \alpha_0 = i_1^*\omega - i_0^*\omega .
\]

If $F \colon M \times \RR \to N$ is smooth with $F \circ i_0 = f$ and $F \circ i_1 = g$, then for closed $\omega$ on $N$,
$g^*\omega - f^*\omega = d(K F^*\omega)$, so $f^* = g^*$ on cohomology. Continuous homotopies may be replaced by smooth ones: continuous maps between manifolds are homotopic to smooth maps, and continuously
homotopic smooth maps are smoothly homotopic \cite[Chapter 2]{hirsch-differential-topology}. A convex $U$ is smoothly contractible, by the straight-line homotopy, so its de Rham cohomology is that of a
point.
\end{proof}

\begin{remark}\label{de-rham:remark-smooth-homotopy}
\emph{Smooth and continuous homotopies.} Here a homotopy between smooth maps $f, g \colon M \to N$ is a smooth map $F \colon M \times \RR \to N$ with $F(\cdot, 0) = f$ and $F(\cdot, 1) = g$, a \emph{smooth
homotopy}, and a smooth homotopy equivalence is defined with smooth maps and smooth homotopies. Nothing is lost: by the Whitney approximation theorem every continuous map between smooth manifolds is
homotopic to a smooth one, and two smooth maps that are homotopic through continuous maps are smoothly homotopic \cite[Chapter 6]{lee-smooth-manifolds}. So homotopic in the sense of
Chapter~\ref{homotopy} and smoothly homotopic agree for smooth maps, and homotopy equivalent manifolds have isomorphic de Rham cohomology, in accordance with the de Rham theorem below, which identifies
de Rham cohomology with singular cohomology, a homotopy invariant.
\end{remark}

\begin{remark}\label{de-rham:remark-k-and-prism}
\emph{The homotopy operator and the prism operator.} Singular homology has the same homotopy invariance, with the same structure of proof: a chain homotopy $P$ with $\partial P + P\partial = g_\# - f_\#$, the
\emph{prism operator} of Theorem~\ref{singular-homology:theorem-homotopy-invariance}, plays the role of $K$ with $dK + Kd = i_1^* - i_0^*$. The two are dual to each other under integration. Let $\sigma \colon \Delta^k \to M$ be a
smooth singular simplex, $F \colon M \times \RR \to N$ a smooth homotopy and $\omega$ a $(k+1)$-form on $N$. The prism $P(\sigma)$ is $F \circ (\sigma \times \id)$ restricted to the $k + 1$ simplices that triangulate
$\Delta^k \times [0, 1]$ (Remark~\ref{singular-homology:remark-prism}), so integrating $\omega$ over it means integrating $(\sigma \times \id)^*F^*\omega = \alpha_t + dt \wedge \beta_t$ over the prism
$\Delta^k \times [0, 1]$. Only the part $dt \wedge \beta_t$ has the right type to contribute, and by Fubini's theorem
\[
\int_{P(\sigma)}\omega = \pm\int_{\Delta^k}\int_0^1\beta_t\,dt = \pm\int_\sigma K(F^*\omega),
\]
with a sign depending only on $k$ and on the orientation conventions. So $K$ is ``integration over the prism, one dimension at a time'': it integrates over the $I$-direction and leaves the integral over
$\Delta^k$ for later. The identities $dK + Kd = i_1^* - i_0^*$ and $\partial P + P\partial = g_\# - f_\#$ correspond to each other under Stokes' theorem, which turns the boundary of the prism into the exterior derivative.
For forms no triangulation is needed; for chains the prism must be cut into simplices, and the cuts cancel in pairs. This compatibility is one ingredient of the de Rham theorem
(Theorem~\ref{de-rham:theorem-de-rham}): the integration map from forms to singular cochains is compatible with the two homotopy invariances.
\end{remark}

The homotopy operator gives explicit primitives. For a star-shaped open set $U \subseteq \RR^n$ (with respect to $0$), let $X = \sum_ix^i\partial_i$ be the position vector
field.

\begin{corollary}[Explicit Poincar\'e lemma]\label{de-rham:corollary-explicit-poincare}
Let $U \subseteq \RR^n$ be star-shaped with respect to $0$ and $\omega$ a closed $k$-form on $U$, $k \geq 1$. Then $\omega = d\eta$ with
\[
\eta_x = \int_0^1t^{k-1}\,(\iota_X\omega)_{tx}\,dt, \qquad\text{that is,}\qquad \eta = \sum_I\Bigl(\int_0^1t^{k-1}\omega_I(tx)\,dt\Bigr)\iota_X(dx^I),
\]
where $\iota_X(dx^{i_1} \wedge \cdots \wedge dx^{i_k}) = \sum_m(-1)^{m-1}x^{i_m}\,dx^{i_1} \wedge \cdots \wedge \widehat{dx^{i_m}} \wedge \cdots \wedge dx^{i_k}$.
\end{corollary}

\begin{proof}
Let $\rho \colon \RR \to [0, 1]$ be smooth and nondecreasing, with $\rho = 0$ on $(-\infty, 0]$ and $\rho = 1$ on $[1, \infty)$, and apply
Theorem~\ref{de-rham:theorem-homotopy-invariance} to $F(x, t) = \rho(t)x$, a smooth map $U \times \RR \to U$ since $U$ is star-shaped. Since
$F^*dx^i = \rho(t)\,dx^i + \rho^{\prime}(t)x^i\,dt$, the terms of $F^*(\omega_I\,dx^{i_1} \wedge \cdots \wedge dx^{i_k})$ containing $dt$ are
$\omega_I(\rho x)\,\rho^{k-1}\rho^{\prime}\sum_m x^{i_m}\,dx^{i_1} \wedge \cdots \wedge dt \wedge \cdots \wedge dx^{i_k}$, with $dt$ in the $m$-th place; moving $dt$ to the front gives the sign
$(-1)^{m-1}$. So $\beta_t = \rho(t)^{k-1}\rho^{\prime}(t)(\iota_X\omega)_{\rho(t)x}$ in the notation of the proof of the theorem, and the substitution $s = \rho(t)$ gives
$K(F^*\omega) = \int_0^1s^{k-1}(\iota_X\omega)_{sx}\,ds = \eta$. As $F(\cdot, 1) = \id$, $F(\cdot, 0)$ is constant, whose pullback kills forms of positive degree, and $d\omega = 0$, the
homotopy formula gives $\omega = d\eta$.
\end{proof}

\begin{example}\label{de-rham:example-explicit-primitives}
\begin{enumerate}
\item For the closed $1$-form $\omega = 2xy\,dx + x^2\,dy$ on $\RR^2$, $\iota_X\omega = 2x^2y + x^2y = 3x^2y$, so $\eta = \int_0^13t^2\,dt \cdot x^2y = x^2y$, and indeed $d(x^2y) = \omega$.
\item For $\omega = dx \wedge dy$ on $\RR^2$, $\iota_X\omega = x\,dy - y\,dx$ and $\eta = \int_0^1t\,dt\,(x\,dy - y\,dx) = \frac12(x\,dy - y\,dx)$, the primitive used in
Example~\ref{differential-forms:example-stokes-checks}(1) to compute areas.
\item For the volume form $dx \wedge dy \wedge dz$ of $\RR^3$ the formula gives $\frac13(x\,dy \wedge dz + y\,dz \wedge dx + z\,dx \wedge dy)$.
\item The slit plane $\RR^2 \setminus \{(x, 0) : x \leq 0\}$ does not contain $0$, but it is star-shaped with respect to $(1, 0)$: a segment from $(1, 0)$ to a point off the
$x$-axis meets the axis only at $(1, 0)$. So the angle form is exact there; a primitive is the polar angle with values in $(-\pi, \pi)$.
\end{enumerate}
\end{example}

\begin{theorem}[Mayer--Vietoris]\label{de-rham:theorem-mayer-vietoris}
Let $M = U \cup V$ with $U$, $V$ open. There is a long exact sequence
\[
\cdots \to H^k_{\mathrm{dR}}(M) \to H^k_{\mathrm{dR}}(U) \oplus H^k_{\mathrm{dR}}(V) \to H^k_{\mathrm{dR}}(U \cap V) \xrightarrow{\ \delta\ } H^{k+1}_{\mathrm{dR}}(M) \to \cdots
\]
\end{theorem}

\begin{proof}
The sequence of complexes $0 \to \Omega^\bullet(M) \to \Omega^\bullet(U) \oplus \Omega^\bullet(V) \to \Omega^\bullet(U \cap V) \to 0$, with maps $\omega \mapsto (\omega|_U, \omega|_V)$ and
$(\alpha, \beta) \mapsto \alpha|_{U \cap V} - \beta|_{U \cap V}$, is exact. Injectivity: a form vanishing on $U$ and on $V$ vanishes. Exactness in the middle: if $\alpha = \beta$ on $U \cap V$, they
glue to a form on $M$. Surjectivity: for a partition of unity $\{\psi_U, \psi_V\}$ subordinate to $\{U, V\}$ (Theorem~\ref{smooth-manifolds:theorem-partitions-of-unity}) and
$\gamma \in \Omega^\bullet(U \cap V)$, the forms $\psi_V\gamma$ on $U$ and $-\psi_U\gamma$ on $V$, both extended by zero (the support of $\psi_V$ does not meet $U \setminus V$, and similarly for
$\psi_U$), have difference $\gamma$ on $U \cap V$. The maps commute with $d$, so Theorem~\ref{homological-algebra:theorem-long-exact-sequence} gives the long exact sequence. Following the
construction of the connecting homomorphism, for a closed $\gamma$ on $U \cap V$ we lift it to $(\psi_V\gamma, -\psi_U\gamma)$, apply $d$, and observe that $d(\psi_V\gamma)$ and
$-d(\psi_U\gamma)$ agree on $U \cap V$, since $d(\psi_U\gamma + \psi_V\gamma) = d\gamma = 0$. So
\[
\delta[\gamma] = [\omega], \qquad \omega = \begin{cases} d(\psi_V\gamma) = d\psi_V \wedge \gamma & \text{on } U,\\ -d(\psi_U\gamma) = -d\psi_U \wedge \gamma & \text{on } V,\end{cases}
\]
a closed form on $M$ supported in $U \cap V$.
\end{proof}

\begin{example}\label{de-rham:example-spheres}
$H^k_{\mathrm{dR}}(S^n) \cong \RR$ for $k = 0, n$ and $0$ otherwise ($n \geq 1$). Let $U = S^n \setminus \{N\}$ and $V = S^n \setminus \{S\}$ for the north and south poles. Both are
diffeomorphic to $\RR^n$ by stereographic projection, so their cohomology is $\RR$ in degree $0$ and $0$ otherwise; and $U \cap V$ is diffeomorphic to $\RR^n \setminus \{0\}$, which is
homotopy equivalent to $S^{n-1}$ (by $x \mapsto x/|x|$, with the inclusion of the unit sphere as homotopy inverse, and the homotopy $(x, t) \mapsto x/((1 - t) + t|x|)$), so
$H^\bullet_{\mathrm{dR}}(U \cap V) \cong H^\bullet_{\mathrm{dR}}(S^{n-1})$ by Theorem~\ref{de-rham:theorem-homotopy-invariance}. The Mayer--Vietoris sequence splits into pieces:
\begin{enumerate}
\item for $k \geq 1$ the terms $H^k(U) \oplus H^k(V)$ vanish, so $\delta \colon H^k_{\mathrm{dR}}(U \cap V) \to H^{k+1}_{\mathrm{dR}}(S^n)$ is an isomorphism, that is
$H^{k+1}(S^n) \cong H^k(S^{n-1})$;
\item in low degrees, $0 \to H^0(S^n) \to H^0(U) \oplus H^0(V) = \RR^2 \to H^0(U \cap V) \to H^1(S^n) \to 0$.
\end{enumerate}
For $n = 1$, $U \cap V$ has two components, so $H^0(U \cap V) = \RR^2$, the map $\RR^2 \to \RR^2$, $(a, b) \mapsto (a - b, a - b)$, has rank $1$, and $H^1(S^1) \cong \RR$, $H^0(S^1) = \RR$. For
$n \geq 2$, $U \cap V$ is connected, the map $\RR^2 \to \RR$ is surjective, and $H^1(S^n) = 0$, $H^0(S^n) = \RR$. Induction on $n$ with (1) gives the result: $H^n(S^n) \cong H^{n-1}(S^{n-1}) \cong \cdots \cong H^1(S^1) \cong \RR$,
and all other groups vanish. For the circle, the generator of $H^1$ is the form $d\theta$, which is closed but not exact; its integral over $S^1$ is $2\pi$
(Exercise~\ref{differential-forms:exercise-pullback-computation}). In general $\int_M \colon H^n_{\mathrm{dR}}(M) \to \RR$ is well defined for compact oriented $M$ by Stokes' theorem
(Theorem~\ref{differential-forms:theorem-stokes}).

The connecting homomorphism can be seen at work on the circle. Take $U$ and $V$ two open arcs covering $S^1$, whose intersection consists of two arcs $A$ and $B$, and let $\gamma$ be the
locally constant function equal to $1$ on $A$ and $0$ on $B$. Then $\delta[\gamma]$ is represented by $d\psi_V$ on $A$ (and $0$ elsewhere), a $1$-form supported in $A$. Near the end of $A$
where $V$ stops, $\psi_V = 0$, as the support of $\psi_V$ lies in $V$; near the end where $U$ stops, $\psi_U = 0$, so $\psi_V = 1$. Hence $\int_Ad\psi_V = \pm1$, and $\delta[\gamma]$ is a nonzero
multiple of the generator of $H^1(S^1)$: a ``bump'' $1$-form concentrated on one arc.
\end{example}

\section{The de Rham theorem}\label{de-rham:section-de-rham-theorem}

\begin{theorem}[de Rham]\label{de-rham:theorem-de-rham}
For every smooth manifold $M$ there are natural isomorphisms
\[
H^k_{\mathrm{dR}}(M) \cong H^k(M, \RR_M) \cong H^k(M; \RR),
\]
the first to sheaf cohomology of the constant sheaf, the second to singular cohomology with real coefficients. Under these isomorphisms the wedge product corresponds to the cup product, and the
isomorphism is induced by integration: the class of a closed $k$-form $\omega$ corresponds to the singular cochain $\sigma \mapsto \int_{\Delta^k}\sigma^*\omega$ on smooth simplices.
\end{theorem}

\begin{proof}
Let $\underline{\Omega}^k$ be the sheaf $U \mapsto \Omega^k(U)$. The sequence of sheaves
\[
0 \to \RR_M \to \underline{\Omega}^0 \xrightarrow{d} \underline{\Omega}^1 \xrightarrow{d} \cdots
\]
is exact: exactness at $\underline{\Omega}^0$ is Proposition~\ref{de-rham:proposition-h0}, and exactness at $\underline{\Omega}^k$ for $k \geq 1$ holds on stalks, because every point has a basis of convex
coordinate neighbourhoods, on which closed forms are exact by the Poincar\'e lemma (Theorem~\ref{de-rham:theorem-homotopy-invariance}). Each $\underline{\Omega}^k$ is a module over the sheaf of smooth
functions, hence fine, hence acyclic on the paracompact space $M$ (Example~\ref{sheaves:example-fine} and
Theorem~\ref{sheaves:theorem-soft-acyclic}). So this is an acyclic resolution of $\RR_M$ and
Proposition~\ref{sheaves:proposition-hypercohomology}(3) gives $H^k(M, \RR_M) \cong H^k(\Gamma(M, \underline{\Omega}^\bullet)) = H^k_{\mathrm{dR}}(M)$. The second isomorphism is
Theorem~\ref{sheaves:theorem-singular-comparison}, since manifolds are locally contractible and paracompact, and it is compatible with products. The description by integration over smooth simplices
is the standard comparison of the two resolutions; see \cite[Chapter 5]{warner-foundations} or \cite[Chapter III]{bredon-sheaf}.
\end{proof}

\begin{corollary}\label{de-rham:corollary-consequences}
For a compact manifold $M$ the spaces $H^k_{\mathrm{dR}}(M)$ are finite-dimensional, they vanish for $k > \dim M$, and $\dim H^k_{\mathrm{dR}}(M) = b_k(M)$, the $k$-th Betti number. For a compact oriented
$n$-manifold, $H^n_{\mathrm{dR}}(M) \cong \RR$ for each connected component, with the isomorphism given by integration.
\end{corollary}

\begin{proof}
Compact manifolds have finitely generated singular homology \cite[Appendix]{hatcher-at}, so Theorem~\ref{de-rham:theorem-de-rham} and the universal coefficient theorem give finite dimension and the
identification with Betti numbers; the vanishing above the dimension holds because there are no nonzero forms of degree $> n$. For the top degree, $H^n(M; \RR) \cong H_0(M; \RR) \cong \RR$ per component by
Poincar\'e duality (Theorem~\ref{poincare-duality:theorem-poincare-duality}), and integration is nonzero on the class of a volume form
(Corollary~\ref{differential-forms:corollary-volume}).
\end{proof}

Explicitly, if $M_1, \ldots, M_c$ are the connected components of the compact oriented manifold $M$ (finitely many, as they are open and cover the compact $M$), then
$\Omega^\bullet(M) = \prod_i\Omega^\bullet(M_i)$, so $H^n_{\mathrm{dR}}(M) = \bigoplus_iH^n_{\mathrm{dR}}(M_i)$, and the isomorphism of the corollary is
\[
H^n_{\mathrm{dR}}(M) \to \RR^c, \qquad [\omega] \mapsto \Bigl(\int_{M_1}\omega, \ldots, \int_{M_c}\omega\Bigr).
\]
So a top degree form on $M$ is exact if and only if its integral over \emph{each} component vanishes; the vanishing of the total integral $\int_M\omega$ is not enough when
$M$ is disconnected. For example, on $M = S^2 \sqcup S^2$ the form equal to the area form on the first sphere and to minus the area form on the second has total integral $0$ but is
not exact.

\begin{proposition}\label{de-rham:proposition-periods}
Let $M$ be a connected manifold and $x_0 \in M$. There is an isomorphism
\[
H^1_{\mathrm{dR}}(M) \xrightarrow{\ \sim\ } \Hom(\pi_1(M, x_0), \RR),
\]
sending $[\omega]$ to the homomorphism whose value on the class of a piecewise smooth loop $\gamma$ is the \emph{period} $\int_\gamma\omega$. In particular $H^1_{\mathrm{dR}}(M) = 0$ if $M$ is simply connected, and a closed
$1$-form is exact if and only if all its periods vanish.
\end{proposition}

\begin{proof}
Compose the isomorphism $H^1_{\mathrm{dR}}(M) \cong H^1(M; \RR)$ of Theorem~\ref{de-rham:theorem-de-rham}, given by integration over smooth simplices, with $H^1(M; \RR) \cong \Hom_\RR(H_1(M; \RR), \RR) = \Hom(H_1(M), \RR)$
(Theorem~\ref{singular-homology:theorem-universal-coefficients} over the field $\RR$) and with $\Hom(H_1(M), \RR) \cong \Hom(\pi_1(M, x_0), \RR)$, which holds by the Hurewicz theorem
(Theorem~\ref{singular-homology:theorem-hurewicz}) because homomorphisms to the abelian group $\RR$ factor through the abelianisation. The composite is the period map.

The last statement can also be seen directly. If $\omega = df$, then $\int_\gamma\omega = f(\gamma(1)) - f(\gamma(0)) = 0$ for loops. Conversely, if all periods vanish, $f(x) = \int_{x_0}^x\omega$ along any piecewise smooth path is
well defined, since two paths differ by a loop, and in a chart around $x$ that is a ball, integrating along segments shows $df = \omega$, by the Poincar\'e lemma on the ball
(Theorem~\ref{de-rham:theorem-homotopy-invariance}): there $\omega = dg$, and $f - g$ is constant.
\end{proof}

\begin{example}\label{de-rham:example-angle-form}
On $M = \RR^2 \setminus \{0\}$ the \emph{angle form} $\omega = \frac{x\,dy - y\,dx}{x^2 + y^2}$ is closed, and its period on the unit circle is $2\pi$, so it is not exact. Since $\pi_1(M) \cong \ZZ$, Proposition~\ref{de-rham:proposition-periods}
gives $H^1_{\mathrm{dR}}(M) = \RR[\omega]$: a closed $1$-form $\eta$ on $M$ is exact if and only if $\int_{S^1}\eta = 0$, and in general $\eta - \frac{1}{2\pi}\big(\int_{S^1}\eta\big)\omega$ is exact. On the simply connected
$\RR^3 \setminus \{0\}$, by contrast, every closed $1$-form is exact.
\end{example}

\begin{counterexample}\label{de-rham:counterexample-infinite-dimensional}
De Rham cohomology of a noncompact manifold can be infinite-dimensional. On $M = \CC \setminus \ZZ$ let $\alpha_n$ be the angle form centred at $n$, the pullback of the angle form of
Counterexample~\ref{differential-forms:counterexample-closed-not-exact} by $z \mapsto z - n$. It is closed, and its integral over the circle $\gamma_m$ of radius $\frac12$ around $m$ is $2\pi$ if $m = n$ and $0$ otherwise (for $m \neq n$ the circle lies in
a half-plane, where $\alpha_n$ is exact). If $\sum_nc_n\alpha_n = d f$ for finitely many nonzero $c_n$, integrating over $\gamma_m$ gives $2\pi c_m = 0$ (Proposition~\ref{differential-forms:proposition-line-integrals}). So the classes
$[\alpha_n]$, $n \in \ZZ$, are linearly independent in $H^1_{\mathrm{dR}}(M)$.
\end{counterexample}

\begin{proposition}\label{de-rham:proposition-finite-quotient}
Let a finite group $G$ act freely on $M$ by diffeomorphisms, and $\pi \colon M \to M/G$ the quotient (Theorem~\ref{smooth-manifolds:theorem-smooth-quotient}). Then $\pi^*$ induces an isomorphism
$H^k_{\mathrm{dR}}(M/G) \cong H^k_{\mathrm{dR}}(M)^G$ onto the classes fixed by $G$.
\end{proposition}

\begin{proof}
As $\pi$ is a local diffeomorphism and $\pi \circ g = \pi$, $\pi^*$ identifies $\Omega^\bullet(M/G)$ with the complex $\Omega^\bullet(M)^G$ of $G$-invariant forms: an invariant form descends, because locally $\pi$ has smooth inverses that differ
by elements of $G$. The averaging operator $A = \frac{1}{|G|}\sum_{g \in G}g^*$ commutes with $d$ and is a projection of $\Omega^\bullet(M)$ onto $\Omega^\bullet(M)^G$. The map $H^k(\Omega^\bullet(M)^G) \to H^k_{\mathrm{dR}}(M)^G$ is injective: if an
invariant $\omega$ equals $d\eta$, then $\omega = A\omega = d(A\eta)$ with $A\eta$ invariant. It is surjective: if $[\omega]$ is $G$-fixed, then $A\omega$ is invariant, closed, and cohomologous to $\omega$, since each $g^*\omega$ is.
\end{proof}

\begin{example}\label{de-rham:example-real-projective}
$H^k_{\mathrm{dR}}(\RR\PP^n) \cong \RR$ for $k = 0$, and for $k = n$ if $n$ is odd; it vanishes otherwise. By Proposition~\ref{de-rham:proposition-finite-quotient} and Example~\ref{de-rham:example-spheres}, only the action of the
antipodal map $a$ on $H^n_{\mathrm{dR}}(S^n) \cong \RR$ matters. The volume form of $S^n$ is $\iota_X\Omega$ restricted to $S^n$, where $\Omega = dx^1 \wedge \cdots \wedge dx^{n+1}$ and $X = \sum x^i\partial_i$; as $a$ is the restriction of $-\mathrm{id}$,
$a^*\Omega = (-1)^{n+1}\Omega$ and $a$ preserves $X$, so $a^*$ acts on $H^n$ by $(-1)^{n+1}$, and the invariant part is $H^n$ for $n$ odd and $0$ for $n$ even. In particular $\RR\PP^n$ is not orientable for $n$ even, by
Corollary~\ref{de-rham:corollary-consequences}.
\end{example}

\begin{theorem}[K\"unneth formula]\label{de-rham:theorem-kunneth}
If $M$ has a finite good cover, the map $\bigoplus_{p+q=k}H^p_{\mathrm{dR}}(M) \otimes H^q_{\mathrm{dR}}(N) \to H^k_{\mathrm{dR}}(M \times N)$, $[\alpha] \otimes [\beta] \mapsto [\mathrm{pr}_M^*\alpha \wedge \mathrm{pr}_N^*\beta]$, is an isomorphism for every
manifold $N$.
\end{theorem}

\begin{proof}
Write $\psi_M$ for the map. It is natural in $M$ for restrictions to open subsets. If $M \cong \RR^m$, then $\mathrm{pr}_N \colon M \times N \to N$ is a homotopy equivalence and $H^\bullet_{\mathrm{dR}}(M) = \RR$ in degree $0$, so $\psi_M$ is
$\mathrm{pr}_N^*$, an isomorphism by Theorem~\ref{de-rham:theorem-homotopy-invariance}. For open $U, V \subseteq M$, tensoring the Mayer--Vietoris sequence of $U, V$ with the vector space $H^q_{\mathrm{dR}}(N)$ and summing over $q$ gives an
exact sequence, which $\psi$ maps to the Mayer--Vietoris sequence of $U \times N$, $V \times N$. The squares involving restrictions commute, and so does the one with the connecting maps: with a partition of unity $\{\rho_U, \rho_V\}$ on
$U \cup V$, the connecting maps are $[\omega] \mapsto [d(\rho_V\omega)]$ on both sides, and for closed $\beta$, $d(\mathrm{pr}_M^*(\rho_V\omega) \wedge \mathrm{pr}_N^*\beta) = \mathrm{pr}_M^*d(\rho_V\omega) \wedge \mathrm{pr}_N^*\beta$. So if $\psi$ is an isomorphism for $U$, $V$ and $U \cap V$, it is one for $U \cup V$
by the five lemma (Lemma~\ref{homological-algebra:lemma-five}). Now let $U_1, \ldots, U_k$ be a finite good cover of $M$; we show by induction on $k$ that $\psi$ is an isomorphism for every open set covered by $k$ open sets forming a good
cover. For $k = 1$ this is the case $\RR^m$. For the induction step apply the five lemma to $W = U_1 \cup \cdots \cup U_{k-1}$ and $U_k$, noting that $W \cap U_k$ is covered by the $k - 1$ sets $U_i \cap U_k$, which form a good cover.
\end{proof}

\begin{corollary}\label{de-rham:corollary-torus}
$H^\bullet_{\mathrm{dR}}(T^n)$ is the exterior algebra on the classes of $dx^1, \ldots, dx^n$, so $b_k(T^n) = \binom nk$; and $H^\bullet_{\mathrm{dR}}(S^1 \times S^2)$ has dimensions $1, 1, 1, 1$.
\end{corollary}

\begin{proof}
By Example~\ref{de-rham:example-spheres}, $H^\bullet(S^1)$ is spanned by $1$ and $[d\theta]$, and $T^n = (S^1)^n$ has a finite good cover; Theorem~\ref{de-rham:theorem-kunneth} and induction give the tensor product of $n$ copies of
$\RR \oplus \RR[d\theta]$, which is the exterior algebra because $d\theta \wedge d\theta = 0$ and the wedge product is graded commutative. The second statement is the same computation with $H^\bullet(S^2)$.
\end{proof}

\begin{remark}\label{de-rham:remark-no-torsion}
\emph{De Rham cohomology is a vector space, and it does not see torsion.} The spaces $\Omega^k(M)$ are real vector spaces and $d$ is $\RR$-linear, so $H^k_{\mathrm{dR}}(M)$ is a real vector space, not just an abelian group.
When it is finite-dimensional, for instance for compact $M$ (Corollary~\ref{de-rham:corollary-consequences}), a choice of basis gives $H^k_{\mathrm{dR}}(M) \cong \RR^{b_k}$, and the only invariant left is the dimension.
A real vector space has no torsion: $mv = 0$ with $m \neq 0$ forces $v = m^{-1}(mv) = 0$. Integer phenomena are therefore lost. By Theorem~\ref{de-rham:theorem-de-rham} and the universal coefficient theorem
(Theorem~\ref{singular-homology:theorem-universal-coefficients}) over the field $\RR$,
\[
H^k_{\mathrm{dR}}(M) \cong H^k(M; \RR) \cong \Hom(H_k(M), \RR),
\]
and $\Hom(\ZZ/m\ZZ, \RR) = 0$, since the image of a generator $x$ would satisfy $m\varphi(x) = 0$. So if $H_k(M) \cong \ZZ^{r} \oplus T$ with $T$ finite, then $b_k = r$: de Rham cohomology records the free part of homology and
forgets $T$. For example (Examples~\ref{de-rham:example-spheres} and~\ref{de-rham:example-real-projective}), $S^3$ and $\RR\PP^3$ have the same de Rham cohomology, $\RR$ in degrees $0$ and $3$, although
$H_1(\RR\PP^3) \cong \ZZ/2\ZZ$ while $H_1(S^3) = 0$; and the Klein bottle has $H_1 \cong \ZZ \oplus \ZZ/2\ZZ$ (Exercise~\ref{singular-homology:exercise-klein-bottle}) and $H^1_{\mathrm{dR}} \cong \RR$
(Example~\ref{de-rham:example-dimension-counting}(2)). To see torsion one needs integer coefficients, singular or \v{C}ech cohomology with values in $\ZZ$. This matters for the Hodge conjecture: its rational form only
concerns $H^{2p}(X; \QQ)$, while torsion classes in $H^{2p}(X; \ZZ)$ are responsible for the first counterexamples to the integral version (Chapter~\ref{known-cases}).
\end{remark}

Because the groups are vector spaces, computations can be done by counting dimensions. The basic tool is the following.

\begin{proposition}\label{de-rham:proposition-dimension-counting}
\begin{enumerate}
\item If $0 \to V_1 \to V_2 \to \cdots \to V_m \to 0$ is an exact sequence of finite-dimensional vector spaces, then $\sum_{i=1}^m(-1)^i\dim V_i = 0$.
\item Let $M = U \cup V$ with $U$, $V$ open. If two of the three manifolds $M$, $U \sqcup V$ and $U \cap V$ have finite-dimensional de Rham cohomology (in all degrees together), then so does the third, and the
\emph{de Rham Euler characteristics} $\chi(W) = \sum_k(-1)^k\dim H^k_{\mathrm{dR}}(W)$ satisfy
\[
\chi(M) = \chi(U) + \chi(V) - \chi(U \cap V).
\]
\item A manifold with a finite good cover, a finite open cover $U_1, \ldots, U_k$ in which every nonempty intersection $U_{i_1} \cap \cdots \cap U_{i_p}$ is diffeomorphic to $\RR^n$, has finite-dimensional de Rham
cohomology.
\end{enumerate}
\end{proposition}

Here $H^k_{\mathrm{dR}}(U \sqcup V) = H^k_{\mathrm{dR}}(U) \oplus H^k_{\mathrm{dR}}(V)$, since a form on a disjoint union is a pair of forms. For compact $M$, $\chi(M)$ is the Euler characteristic of
Chapter~\ref{singular-homology}, because $\dim H^k_{\mathrm{dR}}(M) = b_k(M)$ (Corollary~\ref{de-rham:corollary-consequences}). The formula in (2) is the analogue of the inclusion--exclusion principle
$|U \cup V| = |U| + |V| - |U \cap V|$ for finite sets: the Euler characteristic behaves like a ``number of points'', and for a finite set of points it is exactly that.

\begin{proof}
(1) Put $V_0 = V_{m+1} = 0$ and let $f_i \colon V_i \to V_{i+1}$, $0 \leq i \leq m$, be the maps, so that $f_0 = 0$ and $f_m = 0$. By rank--nullity (Theorem~\ref{linear-algebra:theorem-rank-nullity}), $\dim V_i = \dim\ker f_i + \dim\im f_i$, and exactness says
$\dim\ker f_i = \dim\im f_{i-1}$. So $\sum_i(-1)^i\dim V_i = \sum_i(-1)^i(\dim\im f_{i-1} + \dim\im f_i)$, a telescoping sum equal to $\pm\dim\im f_0 \pm \dim\im f_m = 0$.

(2) First, if $A \xrightarrow{a} B \xrightarrow{b} C$ is exact with $A$ and $C$ finite-dimensional, so is $B$: by rank--nullity $\dim B = \dim\ker b + \dim\im b$, and $\ker b = \im a$ has dimension at most $\dim A$,
while $\im b \subseteq C$. In the Mayer--Vietoris sequence (Theorem~\ref{de-rham:theorem-mayer-vietoris})
\[
\cdots \to H^{k-1}_{\mathrm{dR}}(U \cap V) \to H^k_{\mathrm{dR}}(M) \to H^k_{\mathrm{dR}}(U \sqcup V) \to H^k_{\mathrm{dR}}(U \cap V) \to H^{k+1}_{\mathrm{dR}}(M) \to \cdots
\]
each space sits between a space of each of the other two kinds, so this argument shows that the third kind is finite-dimensional. All the spaces vanish above $\dim M$, so the Mayer--Vietoris sequence is a finite exact
sequence of finite-dimensional spaces; apply (1). In it, $H^k(M)$, $H^k(U) \oplus H^k(V)$ and $H^k(U \cap V)$ occupy the positions $3k + 1$, $3k + 2$, $3k + 3$, with signs $(-1)^{k+1}$, $(-1)^k$, $(-1)^{k+1}$ in
the alternating sum, which gives $-\chi(M) + \chi(U) + \chi(V) - \chi(U \cap V) = 0$.

(3) By induction on the number $k$ of sets in a good cover: for $k = 1$ the manifold is diffeomorphic to $\RR^n$, whose cohomology is $\RR$ in degree $0$ by the Poincar\'e lemma
(Theorem~\ref{de-rham:theorem-homotopy-invariance}). We prove that every open set covered by $k$ sets forming a good cover has finite-dimensional cohomology. For the induction step write the union as $W \cup U_k$ with $W = U_1 \cup \cdots \cup U_{k-1}$; then
$W \cap U_k$ is covered by the (at most) $k - 1$ nonempty sets $U_i \cap U_k$, which form a good cover of it because their intersections are intersections of the $U_j$. By induction $W$ and $W \cap U_k$ have
finite-dimensional cohomology, and so does $U_k \cong \RR^n$; so (2) applies.
\end{proof}

\begin{example}\label{de-rham:example-dimension-counting}
\begin{enumerate}
\item \emph{The torus.} Write $T^2 = S^1 \times S^1$ as the union of two open annuli $U = S^1 \times A$ and $V = S^1 \times B$, where $A$ and $B$ are open arcs covering $S^1$; then $U \cap V$ is a disjoint union of two annuli. Each annulus is
homotopy equivalent to $S^1$ (Theorem~\ref{de-rham:theorem-homotopy-invariance}), which has $\chi = 1 - 1 = 0$ (Example~\ref{de-rham:example-spheres}). So $\chi(T^2) = 0 + 0 - 0 = 0$. Since $T^2$ is compact, connected and orientable, $b_0 = b_2 = 1$ (Proposition~\ref{de-rham:proposition-h0} and
Corollary~\ref{de-rham:corollary-consequences}), and $1 - b_1 + 1 = 0$ gives $b_1 = 2$, as in Corollary~\ref{de-rham:corollary-torus}, but without computing any map in the sequence. In general, the
Euler characteristic of a Mayer--Vietoris decomposition is computed from the pieces alone, while the individual $b_k$ depend on the maps; dimension counting gives the $b_k$ as soon as all but one of them are known.
\item \emph{The Klein bottle.} The Klein bottle is $K = T^2/G$ for the group $G = \{1, \tau\}$ generated by the free involution $\tau(x, y) = (x + \frac12, -y)$ of $T^2 = \RR^2/\ZZ^2$. By Corollary~\ref{de-rham:corollary-torus},
$H^1_{\mathrm{dR}}(T^2)$ has basis $[dx]$, $[dy]$ and $H^2_{\mathrm{dR}}(T^2)$ basis $[dx \wedge dy]$; and $\tau^*dx = dx$, $\tau^*dy = -dy$, $\tau^*(dx \wedge dy) = -dx \wedge dy$. By
Proposition~\ref{de-rham:proposition-finite-quotient} the invariant classes give $H^\bullet_{\mathrm{dR}}(K)$: $\RR$, $\RR[dx]$, $0$ in degrees $0$, $1$, $2$. So $\chi(K) = 1 - 1 + 0 = 0$, as for the torus, while
$H^2_{\mathrm{dR}}(K) = 0$ shows that $K$ is not orientable, since a compact connected oriented surface has $H^2_{\mathrm{dR}} \cong \RR$ (Corollary~\ref{de-rham:corollary-consequences}). Compare with
$H_1(K) \cong \ZZ \oplus \ZZ/2\ZZ$ (Exercise~\ref{singular-homology:exercise-klein-bottle}): the free part $\ZZ$ gives $b_1 = 1$, and the torsion $\ZZ/2\ZZ$ is invisible (Remark~\ref{de-rham:remark-no-torsion}).
\item \emph{Counting with the Euler characteristic.} If a compact connected surface $S$ is the union of open sets whose cohomology is known, (2) of the proposition gives $\chi(S)$, and then
$b_1 = b_0 + b_2 - \chi(S)$ with $b_0 = 1$ and $b_2 = 1$ or $0$ according as $S$ is orientable or not. Exercise~\ref{de-rham:exercise-genus-betti} does this for the surfaces of genus $g$.
\end{enumerate}
\end{example}

\subsection{Cohomology with compact supports}\label{de-rham:subsection-compact-supports}

Poincar\'e duality pairs the cohomology of a manifold with its cohomology with compact supports (Definition~\ref{de-rham:definition-de-rham}). For compact manifolds the two coincide; for $\RR^n$ they are as different
as possible: $H^\bullet_{\mathrm{dR}}(\RR^n)$ is concentrated in degree $0$, and $H^\bullet_{\mathrm{dR},c}(\RR^n)$, as we now show, in degree $n$. We first compute the latter, then prove a Mayer--Vietoris sequence for it,
and deduce duality by induction on a good cover.

\begin{lemma}[Integration along the fibre]\label{de-rham:lemma-integration-fibre}
On $\RR^{m+1} = \RR^m \times \RR$ with coordinates $(x, t)$ write every compactly supported $k$-form uniquely as $\omega = \sum_If_I\,dx^I + \sum_Jg_J\,dx^J \wedge dt$ ($|I| = k$, $|J| = k - 1$). Fix $e \in C^\infty_c(\RR)$ with
$\int_\RR e = 1$ and let $A(t) = \int_{-\infty}^te$. Define
\[
\begin{aligned}
\pi_*\omega &= \sum_J\Bigl(\int_\RR g_J(x, t)\,dt\Bigr)dx^J, \qquad e_*\varphi = \varphi \wedge e(t)\,dt, \\
K\omega &= \sum_J\Bigl(\int_{-\infty}^tg_J(x, s)\,ds - A(t)\int_\RR g_J(x, s)\,ds\Bigr)dx^J,
\end{aligned}
\]
for $\omega \in \Omega^k_c(\RR^{m+1})$ and $\varphi \in \Omega^{k-1}_c(\RR^m)$ (pulled back to $\RR^{m+1}$). Then $\pi_* \colon \Omega^k_c(\RR^{m+1}) \to \Omega^{k-1}_c(\RR^m)$, $e_*$ and $K$ preserve compact supports,
$\pi_*d = d\pi_*$, $e_*d = de_*$, $\pi_*e_* = \id$, and
\[
\omega - e_*\pi_*\omega = (-1)^{k-1}(dK - Kd)\omega \qquad (\omega \in \Omega^k_c(\RR^{m+1})).
\]
Hence $\pi_*$ and $e_*$ induce inverse isomorphisms $H^k_{\mathrm{dR},c}(\RR^{m+1}) \cong H^{k-1}_{\mathrm{dR},c}(\RR^m)$, and $\int_{\RR^m}\pi_*\omega = \int_{\RR^{m+1}}\omega$ for $k = m + 1$ (with the orientation $dx^1 \wedge \cdots \wedge dx^m \wedge dt$).
\end{lemma}

\begin{proof}
\emph{Step 1: supports.} If $\omega$ vanishes outside $C \times [-T, T]$ with $C$ compact, then $\pi_*\omega$ vanishes outside $C$, and $K\omega$ outside $C \times [-T^{\prime}, T^{\prime}]$ where $\operatorname{supp}e \subseteq [-T^{\prime}, T^{\prime}]$ and
$T^{\prime} \geq T$: for $t < -T^{\prime}$ both terms vanish, and for $t > T^{\prime}$ they are both $\int_\RR g_J$.

\emph{Step 2: commutation with $d$.} For $\omega = f\,dx^I$, $d\omega = \sum_i\partial_if\,dx^i \wedge dx^I + (-1)^k\partial_tf\,dx^I \wedge dt$, so $\pi_*d\omega = (-1)^k(\int\partial_tf\,dt)dx^I = 0 = d\pi_*\omega$, since $f$ has compact
support in $t$. For $\omega = g\,dx^J \wedge dt$, $d\omega = \sum_i\partial_ig\,dx^i \wedge dx^J \wedge dt$ and $\pi_*d\omega = \sum_i(\int\partial_ig\,dt)dx^i \wedge dx^J = d\pi_*\omega$, differentiating under the integral sign.
Next, $d(e_*\varphi) = d\varphi \wedge e\,dt \pm \varphi \wedge e^{\prime}\,dt \wedge dt = e_*d\varphi$, and $\pi_*e_*\varphi = \varphi\int e = \varphi$.

\emph{Step 3: the homotopy.} For $\omega = f\,dx^I$: $K\omega = 0$, $e_*\pi_*\omega = 0$ and $Kd\omega = (-1)^k(\int_{-\infty}^t\partial_sf\,ds - A(t)\int_\RR\partial_sf\,ds)dx^I = (-1)^kf\,dx^I$, so
$(dK - Kd)\omega = (-1)^{k-1}\omega$. For $\omega = g\,dx^J \wedge dt$ let $G = \int_{-\infty}^tg\,ds - A(t)\bar g$ with $\bar g(x) = \int_\RR g(x, s)\,ds$, so $K\omega = G\,dx^J$, $\partial_tG = g - e\bar g$ and
$\partial_iG = \int_{-\infty}^t\partial_ig\,ds - A(t)\int_\RR\partial_ig\,ds$. Then
\[
\begin{aligned}
dK\omega &= \sum_i\partial_iG\,dx^i \wedge dx^J + (-1)^{k-1}(g - e\bar g)\,dx^J \wedge dt, \\
Kd\omega &= \sum_i\partial_iG\,dx^i \wedge dx^J,
\end{aligned}
\]
so $(dK - Kd)\omega = (-1)^{k-1}(\omega - \bar g\,dx^J \wedge e\,dt) = (-1)^{k-1}(\omega - e_*\pi_*\omega)$. By linearity the identity holds for all $\omega$.

\emph{Step 4: conclusion.} By Step 2, $\pi_*$ and $e_*$ are chain maps; $\pi_*e_* = \id$, and by Step 3 $e_*\pi_*$ induces the identity on cohomology, since for closed $\omega$ the difference $\omega - e_*\pi_*\omega = (-1)^{k-1}dK\omega$
is the differential of a compactly supported form. The statement on integrals is Fubini's theorem (Theorem~\ref{measure-theory:theorem-fubini}) applied to $\int g\,dx\,dt$.
\end{proof}

\begin{proposition}[Poincar\'e lemma with compact supports]\label{de-rham:proposition-compact-poincare}
$H^k_{\mathrm{dR},c}(\RR^n) = 0$ for $k \neq n$, and $\int \colon H^n_{\mathrm{dR},c}(\RR^n) \to \RR$ is an isomorphism: a compactly supported $n$-form on $\RR^n$ is $d$ of a compactly supported form if and only if its integral is $0$.
\end{proposition}

\begin{proof}
By induction on $n$ with Lemma~\ref{de-rham:lemma-integration-fibre}: $H^k_{\mathrm{dR},c}(\RR^n) \cong H^{k-1}_{\mathrm{dR},c}(\RR^{n-1}) \cong \cdots \cong H^{k-n}_{\mathrm{dR},c}(\RR^0)$, which is $\RR$ for $k = n$ and $0$ otherwise
($\RR^0$ is a point, whose only forms are the constants in degree $0$). The isomorphism in degree $n$ is compatible with integration at each step, so it is $\int$.
\end{proof}

The contrast with $H^\bullet_{\mathrm{dR}}(\RR^n)$ is explained by Stokes' theorem: a compactly supported $n$-form with nonzero integral cannot be $d$ of a compactly supported form, although it is $d$ of some form. Compactly
supported cohomology is not a homotopy invariant: $\RR^n$ and a point have different $H^\bullet_{\mathrm{dR},c}$. It is functorial for proper maps and, covariantly, for inclusions of open sets.

\begin{proposition}[Mayer--Vietoris with compact supports]\label{de-rham:proposition-mayer-vietoris-compact}
Let $M = U \cup V$ with $U$, $V$ open, and $\{\psi_U, \psi_V\}$ a partition of unity subordinate to $\{U, V\}$. There is a long exact sequence
\[
\cdots \to H^k_{\mathrm{dR},c}(U \cap V) \to H^k_{\mathrm{dR},c}(U) \oplus H^k_{\mathrm{dR},c}(V) \to H^k_{\mathrm{dR},c}(M) \xrightarrow{\ \delta_c\ } H^{k+1}_{\mathrm{dR},c}(U \cap V) \to \cdots,
\]
where the first map is $\gamma \mapsto (\gamma, -\gamma)$ and the second $(\alpha, \beta) \mapsto \alpha + \beta$ (all forms extended by zero), and $\delta_c[\tau] = [d\psi_U \wedge \tau]$.
\end{proposition}

\begin{proof}
The sequence $0 \to \Omega^\bullet_c(U \cap V) \to \Omega^\bullet_c(U) \oplus \Omega^\bullet_c(V) \to \Omega^\bullet_c(M) \to 0$ of complexes is exact. Injectivity is clear. In the middle, if $\alpha + \beta = 0$ then $\alpha = -\beta$ has support
contained in $\operatorname{supp}\alpha \cap \operatorname{supp}\beta$, a compact subset of $U \cap V$. On the right, $\tau = \psi_U\tau + \psi_V\tau$, and $\operatorname{supp}(\psi_U\tau) \subseteq \operatorname{supp}\psi_U \cap \operatorname{supp}\tau$
is compact (closed in the compact $\operatorname{supp}\tau$) and contained in $U$. Theorem~\ref{homological-algebra:theorem-long-exact-sequence} gives the long exact sequence. For the connecting map, lift a closed $\tau$ to
$(\psi_U\tau, \psi_V\tau)$ and apply $d$: $(d\psi_U \wedge \tau, d\psi_V \wedge \tau)$, which is the image of $d\psi_U \wedge \tau$, a form with compact support in $U \cap V$ since $d\psi_U = -d\psi_V$ vanishes outside
$\operatorname{supp}\psi_U \cap \operatorname{supp}\psi_V$.
\end{proof}

\begin{theorem}[Poincar\'e duality in de Rham cohomology]\label{de-rham:theorem-poincare-duality-de-rham}
Let $M$ be an oriented $n$-manifold. The pairing
\[
H^k_{\mathrm{dR}}(M) \times H^{n-k}_{\mathrm{dR},c}(M) \to \RR, \qquad ([\alpha], [\beta]) \mapsto \int_M\alpha \wedge \beta,
\]
is well defined and nondegenerate if $M$ has a finite good cover (for instance if $M$ is compact); in that case $H^k_{\mathrm{dR}}(M) \cong H^{n-k}_{\mathrm{dR},c}(M)^*$. For compact oriented $M$ this is the
statement that $\int_M\alpha \wedge \beta$ is a perfect pairing between $H^k_{\mathrm{dR}}(M)$ and $H^{n-k}_{\mathrm{dR}}(M)$.
\end{theorem}

\begin{proof}
\emph{Step 1: the pairing is well defined.} If $\alpha$ is closed and $\beta = d\gamma$ with $\gamma$ compactly supported, then $\alpha \wedge \beta = \pm d(\alpha \wedge \gamma)$ has compact support and integrates to zero by Stokes' theorem
(Theorem~\ref{differential-forms:theorem-stokes}); similarly if $\alpha = d\eta$ and $\beta$ is closed. Let $\mathrm{PD}_M \colon H^k_{\mathrm{dR}}(M) \to H^{n-k}_{\mathrm{dR},c}(M)^*$ be the induced map; we show that it is an isomorphism.

\emph{Step 2: $M = \RR^n$.} $H^k_{\mathrm{dR}}(\RR^n)$ is $\RR$ (the constants) for $k = 0$ and $0$ otherwise (Theorem~\ref{de-rham:theorem-homotopy-invariance}), and $H^{n-k}_{\mathrm{dR},c}(\RR^n)$ is $\RR$, by integration, for $k = 0$ and
$0$ otherwise (Proposition~\ref{de-rham:proposition-compact-poincare}). The pairing of the constant $c$ with $[\beta]$ is $c\int\beta$, which is nondegenerate.

\emph{Step 3: Mayer--Vietoris.} Let $M = U \cup V$. Dualising the sequence of Proposition~\ref{de-rham:proposition-mayer-vietoris-compact} (dualising preserves exactness of sequences of vector spaces) and placing it under the sequence of
Theorem~\ref{de-rham:theorem-mayer-vietoris}, with the maps $\mathrm{PD}$ as vertical arrows, gives a ladder of exact sequences. Its squares commute up to sign. For the restriction square, $\int_M\omega \wedge (\alpha + \beta) = \int_U\omega \wedge \alpha + \int_V\omega \wedge \beta$;
for the difference square, $\int_{U \cap V}(\alpha - \beta) \wedge \gamma = \int_U\alpha \wedge \gamma + \int_V\beta \wedge (-\gamma)$; and for the connecting maps, with $\gamma$ closed of degree $k$ on $U \cap V$ and $\tau$ closed with compact support on $M$,
using $\delta[\gamma] = [d\psi_V \wedge \gamma]$ and $d\psi_U = -d\psi_V$,
\[
\int_{U \cap V}\gamma \wedge d\psi_U \wedge \tau = (-1)^k\int_{U \cap V}d\psi_U \wedge \gamma \wedge \tau = (-1)^{k+1}\int_Md\psi_V \wedge \gamma \wedge \tau .
\]
Changing the sign of some vertical maps makes the ladder commute without changing which of them are isomorphisms. By the five lemma (Lemma~\ref{homological-algebra:lemma-five}), if $\mathrm{PD}$ is an isomorphism for $U$, $V$ and $U \cap V$ in all degrees,
it is one for $M$.

\emph{Step 4: induction on a good cover.} Suppose $M$ has a good cover $\{U_1, \ldots, U_p\}$: all nonempty finite intersections are diffeomorphic to $\RR^n$. We show by induction on $p$ that $\mathrm{PD}$ is an isomorphism for every such
$M$. For $p = 1$ this is Step 2. In general let $U = U_1 \cup \cdots \cup U_{p-1}$ and $V = U_p$; then $U$ and $U \cap V = \bigcup_{i<p}(U_i \cap U_p)$ have good covers with $p - 1$ sets, and $V \cong \RR^n$, so Step 3 and the induction
hypothesis conclude. Compact manifolds have finite good covers \cite[Chapter I, Section 5]{bott-tu}; for them $\Omega_c = \Omega$, which gives the last statement.
\end{proof}

Alternatively, for compact $M$, the theorem follows from Theorem~\ref{de-rham:theorem-de-rham} and topological Poincar\'e duality (Corollary~\ref{poincare-duality:corollary-cup-pairing}), since the wedge product corresponds to the cup product.

\begin{theorem}[Cartan--Chevalley--Eilenberg]\label{de-rham:theorem-lie-group-cohomology}
Let $G$ be a compact connected Lie group. Every de Rham class of $G$ has a unique representative that is invariant under left and right translations, and the bi-invariant forms are closed. Consequently
$H^\bullet_{\mathrm{dR}}(G) \cong (\Lambda^\bullet\mathfrak{g}^*)^{G}$, the space of $\Ad(G)$-invariant alternating forms on the Lie algebra.
\end{theorem}

\begin{proof}
$G \times G$ acts on $G$ by $(a, b) \cdot x = axb^{-1}$; let $\mu$ be the Haar probability measure on the compact connected group $K = G \times G$ (Theorem~\ref{lie-groups:theorem-haar}) and, for a form $\omega$, let
$A\omega = \int_Kk^*\omega\,d\mu(k)$, defined pointwise; $A\omega$ is smooth and bi-invariant, $A$ commutes with $d$ (differentiation under the integral sign), and $A\omega = \omega$ for bi-invariant $\omega$.

\emph{Averaging does not change classes.} Let $\omega$ be closed. Each $k \in K$ is connected to the identity by a smooth path, so $k^*\omega$ is cohomologous to $\omega$ (Theorem~\ref{de-rham:theorem-homotopy-invariance}). For a closed form $\eta$
of complementary degree, Fubini's theorem gives $\int_GA\omega \wedge \eta = \int_K(\int_Gk^*\omega \wedge \eta)\,d\mu(k) = \int_G\omega \wedge \eta$, since the inner integral depends only on the class of $k^*\omega$ (Stokes). By Poincar\'e duality
(Theorem~\ref{de-rham:theorem-poincare-duality-de-rham}), $[A\omega] = [\omega]$.

\emph{Bi-invariant forms are closed.} The inversion $\iota(x) = x^{-1}$ satisfies $\iota \circ L_a = R_{a^{-1}} \circ \iota$, so it maps bi-invariant forms to bi-invariant forms, and $d_e\iota = -\mathrm{id}$ (differentiate $\exp(tX)^{-1} = \exp(-tX)$).
A bi-invariant form is determined by its value at $e$, so $\iota^*\omega = (-1)^k\omega$ for a bi-invariant $k$-form. As $d\omega$ is bi-invariant of degree $k + 1$, $(-1)^{k+1}d\omega = \iota^*d\omega = d\iota^*\omega = (-1)^kd\omega$, so $d\omega = 0$.

\emph{Uniqueness.} If a bi-invariant $\omega$ is exact, $\omega = d\eta$, then $\omega = A\omega = d(A\eta) = 0$, because $A\eta$ is bi-invariant, hence closed. So bi-invariant forms map isomorphically onto $H^\bullet_{\mathrm{dR}}(G)$. A
left-invariant form is determined by its value $\omega_e \in \Lambda^\bullet\mathfrak{g}^*$, and it is also right-invariant if and only if $\omega_e$ is invariant under $\Ad(g) = d_e(L_g \circ R_{g^{-1}})$.
\end{proof}

\begin{example}\label{de-rham:example-su2}
For $G = \mathrm{SU}(2)$, $\mathfrak{g} = \mathfrak{su}(2) \cong \RR^3$ and $\Ad$ maps $\mathrm{SU}(2)$ onto $\mathrm{SO}(3)$ (Exercise~\ref{lie-groups:exercise-su2}). The $\mathrm{SO}(3)$-invariants in $\Lambda^\bullet\RR^3$ are the constants and the
multiples of the determinant $e_1^* \wedge e_2^* \wedge e_3^*$, since $\Lambda^1$ and $\Lambda^2 \cong \RR^3$ are irreducible nontrivial representations. So $H^\bullet(\mathrm{SU}(2))$ is $\RR$ in degrees $0$ and $3$, in accordance with
$\mathrm{SU}(2) \cong S^3$.
\end{example}

\section{Degree theory}\label{de-rham:section-degree}

A smooth map between compact oriented manifolds of the same dimension has an integer, its degree, which counts with signs how often it covers a point. With forms the degree appears as the factor by which the map
multiplies integrals of top-degree forms; the two descriptions agree, and together they give quick proofs of classical theorems (the hairy ball theorem, Brouwer's fixed point theorem, the fundamental theorem of
algebra). The topological degree of Chapter~\ref{singular-homology} is the same invariant, defined with singular homology.

\begin{lemma}\label{de-rham:lemma-top-compact}
Let $M$ be a connected oriented $n$-manifold without boundary. Then $\int_M \colon H^n_{\mathrm{dR},c}(M) \to \RR$ is an isomorphism: a compactly supported $n$-form is $d$ of a compactly supported form if and only if its integral vanishes.
\end{lemma}

\begin{proof}
\emph{Step 1: well defined and surjective.} Exact compactly supported forms integrate to zero by Stokes' theorem, and a bump form in a chart has integral $1$.

\emph{Step 2: forms in one ball.} Call an open set a \emph{ball} if it is the domain of an orientation-preserving chart onto $\RR^n$. If $\omega$ is supported in a ball $B$ and $\int\omega = 0$, then $\omega = d\eta$ with $\eta$ supported in $B$, by
Proposition~\ref{de-rham:proposition-compact-poincare}; extending $\eta$ by zero, $\omega$ is exact in $\Omega^\bullet_c(M)$.

\emph{Step 3: moving between balls.} Fix a ball $B_0$ and $\omega_0 \in \Omega^n_c(B_0)$ with $\int\omega_0 = 1$. Let $S$ be the set of points $p$ having a ball neighbourhood $B$ such that some (hence, by Step 2, every)
$\omega_B \in \Omega^n_c(B)$ with $\int\omega_B = 1$ is cohomologous to $\omega_0$ in $H^n_{\mathrm{dR},c}(M)$. If $B$ and $B^{\prime}$ are balls that meet, a form $\omega^{\prime\prime}$ with integral $1$ supported in a ball inside $B \cap B^{\prime}$ is
cohomologous both to $\omega_B$ and to $\omega_{B^{\prime}}$ (Step 2 in $B$ and in $B^{\prime}$). Hence $S$ is open, its complement is open, and $S \ni p_0$ for $p_0 \in B_0$; as $M$ is connected, $S = M$.

\emph{Step 4: general forms.} Cover the support of $\omega$ by finitely many balls $B_1, \ldots, B_r$ and write $\omega = \sum_i\psi_i\omega$ with a partition of unity subordinate to $\{B_1, \ldots, B_r, M \setminus \operatorname{supp}\omega\}$
(Theorem~\ref{smooth-manifolds:theorem-partitions-of-unity}). By Step 3, $\psi_i\omega$ is cohomologous to $(\int\psi_i\omega)\omega_0$, so $\omega$ is cohomologous to $(\int\omega)\omega_0$, which is exact if $\int\omega = 0$.
\end{proof}

\begin{definition}\label{de-rham:definition-degree}
Let $f \colon M \to N$ be a proper smooth map between oriented $n$-manifolds without boundary, $N$ connected and $M$ nonempty. The \emph{degree} of $f$ is the number $\deg f$ such that
\[
\int_Mf^*\omega = \deg f\int_N\omega \qquad \text{for all } \omega \in \Omega^n_c(N).
\]
\end{definition}

The degree exists: by Lemma~\ref{de-rham:lemma-top-compact} the linear form $[\omega] \mapsto \int_Mf^*\omega$ on $H^n_{\mathrm{dR},c}(N) \cong \RR$ is a multiple of $\int_N$ (it is well defined because $f^*$ maps exact
compactly supported forms to exact compactly supported forms, $f$ being proper). It is not obvious from the definition that the degree is an integer; this follows from the local formula.

\begin{theorem}\label{de-rham:theorem-degree-local}
Let $f \colon M \to N$ be as in Definition~\ref{de-rham:definition-degree} and $y \in N$ a regular value of $f$. Then $f^{-1}(y)$ is finite and
\[
\deg f = \sum_{x \in f^{-1}(y)}\operatorname{sign}\det(d_xf),
\]
where the sign is $+1$ if $d_xf$ preserves orientation and $-1$ otherwise. In particular $\deg f \in \ZZ$, and $\deg f = 0$ if $f$ is not surjective.
\end{theorem}

\begin{proof}
\emph{Step 1: the fibre.} By the inverse function theorem every $x \in f^{-1}(y)$ has a neighbourhood mapped diffeomorphically onto a neighbourhood of $y$; so $f^{-1}(y)$ is discrete, and it is compact as $f$ is proper. Hence it is finite,
$\{x_1, \ldots, x_r\}$ (possibly empty).

\emph{Step 2: an evenly covered neighbourhood.} Choose disjoint open $U_i \ni x_i$ mapped diffeomorphically by $f$ onto open sets $V_i \ni y$. The set $C = f(M \setminus \bigcup_iU_i)$ is closed, since proper maps to manifolds are closed, and does not
contain $y$. Let $V$ be a connected neighbourhood of $y$ inside $\bigcap_iV_i \setminus C$ (or inside $N \setminus C$ if $r = 0$). Then $f^{-1}(V) = \bigsqcup_i(U_i \cap f^{-1}(V))$ and $f$ maps each piece diffeomorphically onto $V$.

\emph{Step 3: computation.} Take $\omega \in \Omega^n_c(V)$ with $\int_N\omega = 1$. Then $f^*\omega$ is supported in $f^{-1}(V)$ and $\int_Mf^*\omega = \sum_i\int_{U_i \cap f^{-1}(V)}f^*\omega$. Since $V$ and the pieces are connected, $f$ either
preserves or reverses orientation on each piece, according to the sign of $\det d_{x_i}f$, and by Proposition~\ref{differential-forms:proposition-integral-parametrisation} each integral is $\pm\int_V\omega = \pm1$. If $f$ is not surjective, a point
outside the image is a regular value with empty fibre.
\end{proof}

By Sard's theorem (Theorem~\ref{smooth-manifolds:theorem-sard}) regular values exist, so the formula always computes the degree; in particular the signed count does not depend on the regular value chosen, which is not evident
directly.

\begin{proposition}\label{de-rham:proposition-degree-properties}
\begin{enumerate}
\item $\deg(g \circ f) = \deg g \cdot \deg f$, and a diffeomorphism has degree $\pm1$ according as it preserves or reverses orientation.
\item If $H \colon M \times [0, 1] \to N$ is smooth and proper, then $\deg H_0 = \deg H_1$, where $H_t = H(\cdot, t)$. In particular homotopic maps between compact oriented manifolds have the same degree.
\end{enumerate}
\end{proposition}

\begin{proof}
(1) $\int(g \circ f)^*\omega = \int f^*(g^*\omega) = \deg f\int g^*\omega = \deg f\deg g\int\omega$; for a diffeomorphism use Theorem~\ref{de-rham:theorem-degree-local} at any point.

(2) For $\omega \in \Omega^n_c(N)$, $H^*\omega$ is a compactly supported $n$-form on the manifold with boundary $M \times [0, 1]$, closed as $d\omega = 0$. By Stokes' theorem,
$0 = \int_{M \times [0, 1]}dH^*\omega = \pm\bigl(\int_MH_1^*\omega - \int_MH_0^*\omega\bigr)$, since $\partial(M \times [0, 1]) = M \times \{1\} \cup M \times \{0\}$ with opposite induced orientations.
\end{proof}

\begin{example}\label{de-rham:example-degrees}
\begin{enumerate}
\item The map $z \mapsto z^k$ of $S^1 \subseteq \CC$ has degree $k$: the regular value $1$ has $|k|$ preimages, at which the map preserves orientation if $k > 0$ and reverses it if $k < 0$; for $k = 0$ the map is constant.
\item The antipodal map $a(x) = -x$ of $S^n$ has degree $(-1)^{n+1}$. The form $\sigma = \sum_{i=0}^n(-1)^ix_i\,dx_0 \wedge \cdots \widehat{dx_i} \cdots \wedge dx_n$ restricts to a volume form on $S^n$ for the outward orientation, and
$a^*\sigma = (-1)^{n+1}\sigma$ since $\sigma$ is homogeneous of degree $n + 1$ in the coordinates.
\item A polynomial $p(z) = z^d + c_{d-1}z^{d-1} + \cdots + c_0$, $d \geq 1$, is a proper map $\CC \to \CC$ of degree $d$. It is proper since $|p(z)| \geq |z|^d/2$ for $|z| \geq R$, with $R$ depending only on the $c_i$. The homotopy
$p_t(z) = z^d + t(c_{d-1}z^{d-1} + \cdots + c_0)$, $t \in [0, 1]$, satisfies the same bound uniformly in $t$, so it is proper, and $\deg p = \deg z^d$ by Proposition~\ref{de-rham:proposition-degree-properties}(2). For $z^d$, the regular value $1$ has the $d$
preimages $e^{2\pi ij/d}$, and a holomorphic map preserves orientation where its derivative is nonzero (its real Jacobian is $|p^{\prime}(z)|^2 > 0$). So $\deg p = d \neq 0$, and $p$ is surjective: \emph{every nonconstant complex polynomial has a root}.
\end{enumerate}
\end{example}

\begin{remark}\label{de-rham:remark-degree-applications}
The applications of the degree in Corollary~\ref{singular-homology:corollary-sphere-applications} have short proofs with forms in the smooth setting. \emph{No smooth retraction:} if $r \colon \bar B^n \to S^{n-1}$ were smooth with $r(x) = x$ on $S^{n-1}$, and
$\sigma$ is a volume form on $S^{n-1}$, then by Stokes' theorem
\[
0 < \int_{S^{n-1}}\sigma = \int_{S^{n-1}}r^*\sigma = \int_{\bar B^n}d(r^*\sigma) = \int_{\bar B^n}r^*(d\sigma) = 0,
\]
because $d\sigma$ is an $n$-form on the $(n-1)$-manifold $S^{n-1}$. The ray construction of Corollary~\ref{homotopy:corollary-brouwer} turns a smooth self-map of $\bar B^n$ without fixed points into a smooth retraction, and a continuous self-map
without fixed points can be replaced by a smooth one by uniform approximation (Theorem~\ref{real-analysis:theorem-stone-weierstrass}), which gives Brouwer's theorem. \emph{Hairy ball:} a nowhere vanishing unit tangent field $v$ on $S^n$
gives the smooth homotopy $\cos(\pi t)x + \sin(\pi t)v(x)$ from the identity to the antipodal map, so $1 = (-1)^{n+1}$ by Proposition~\ref{de-rham:proposition-degree-properties}(2) and Example~\ref{de-rham:example-degrees}(2).
\end{remark}

\begin{remark}\label{de-rham:remark-degree-comparison}
For a smooth map $f \colon S^n \to S^n$ the degree of Definition~\ref{de-rham:definition-degree} agrees with the homological degree of Definition~\ref{singular-homology:definition-degree}. Indeed, by Exercise~\ref{de-rham:exercise-degree-de-rham} the map $f^*$ on $H^n_{\mathrm{dR}}(S^n) \cong \RR$ is multiplication by the first degree. By the naturality of the de Rham isomorphism (Theorem~\ref{de-rham:theorem-de-rham}), $f^*$ corresponds to the dual of $f_* \colon H_n(S^n; \RR) \to H_n(S^n; \RR)$, which is multiplication by the second degree, since $H_n(S^n; \RR) = H_n(S^n) \otimes \RR$. So the two degrees agree; in particular the degree of a smooth map, defined with integrals, is an integer for this second reason too.
\end{remark}

\begin{counterexample}\label{de-rham:counterexample-degree}
The hypotheses of the degree theory are needed. Properness: the inclusion of the open interval $(0, 1)$ into $\RR$ is a local diffeomorphism, and the regular value $\frac12$ has one preimage while $2$ has none; there is no well
defined degree, and indeed the inclusion does not pull back compactly supported forms to compactly supported forms. Orientability: on $\RR\PP^2$ there are no orientations to count signs with, and the double cover
$S^2 \to \RR\PP^2$ has a well-defined count only modulo $2$ (here $0$). Compactness of the ball in Brouwer's theorem: the translation $x \mapsto x + 1$ of $\RR$, and the map $x \mapsto \frac{x+1}{2}$ of the open interval $(0, 1)$, have no fixed
points.
\end{counterexample}

\section{Exercises}\label{de-rham:section-exercises}

\begin{exercise}\label{de-rham:exercise-torus}
Compute $H^\bullet_{\mathrm{dR}}(T^n)$ and show that it is the exterior algebra on $H^1$, generated by the forms $d\theta^1, \ldots, d\theta^n$.
\end{exercise}

\begin{solution}
By Theorem~\ref{de-rham:theorem-de-rham} and Example~\ref{singular-homology:example-torus}, $H^\bullet_{\mathrm{dR}}(T^n) \cong \Lambda^\bullet\RR^n$. The classes of the closed forms $d\theta^i$, which are
not exact since $\int_{S^1_i}d\theta^i = 2\pi$, generate: their products are the $\binom{n}{k}$ linearly independent classes of $d\theta^{i_1} \wedge \cdots \wedge d\theta^{i_k}$, which can be detected by
integration over the corresponding subtori.
\end{solution}

\begin{exercise}\label{de-rham:exercise-compactly-supported-line}
Show that $H^0_{\mathrm{dR},c}(\RR) = 0$ and $H^1_{\mathrm{dR},c}(\RR) \cong \RR$, with the isomorphism given by integration; so compactly supported de Rham cohomology is not homotopy invariant.
\end{exercise}

\begin{solution}
A compactly supported function with $df = 0$ is constant, hence zero. The map $\int_\RR \colon \Omega^1_c(\RR) \to \RR$ is surjective, and $f\,dx$ integrates to zero if and only if $F(x) = \int_{-\infty}^xf$ has
compact support, in which case $f\,dx = dF$ is exact.
\end{solution}

\begin{exercise}\label{de-rham:exercise-degree-de-rham}
Let $f \colon M \to N$ be a smooth map of compact connected oriented $n$-manifolds. Show that $f^* \colon H^n_{\mathrm{dR}}(N) \to H^n_{\mathrm{dR}}(M)$ is multiplication by the degree of $f$, and that a map of
nonzero degree is surjective.
\end{exercise}

\begin{solution}
On compact manifolds every form has compact support and $f$ is proper, so $H^n_{\mathrm{dR}} = H^n_{\mathrm{dR},c}$ and $\int_N$ is an isomorphism $H^n_{\mathrm{dR}}(N) \to \RR$, likewise for $M$ (Lemma~\ref{de-rham:lemma-top-compact}). By
Definition~\ref{de-rham:definition-degree}, $\int_Mf^*\omega = \deg f\int_N\omega$, so under these identifications $f^*$ is multiplication by $\deg f$. Surjectivity of maps of nonzero degree is the last part of
Theorem~\ref{de-rham:theorem-degree-local}.
\end{solution}

\begin{exercise}\label{de-rham:exercise-kunneth-de-rham}
Prove the K\"unneth formula $H^\bullet_{\mathrm{dR}}(M \times N) \cong H^\bullet_{\mathrm{dR}}(M) \otimes H^\bullet_{\mathrm{dR}}(N)$ for manifolds with finite good covers, by induction with Mayer--Vietoris.
\end{exercise}

\begin{exercise}\label{de-rham:exercise-punctured-plane}
Let $M = \RR^2 \setminus \{p_1, \ldots, p_k\}$ for distinct points $p_i$. Show that $H^1_{\mathrm{dR}}(M) \cong \RR^k$, with a basis given by the angle forms $\alpha_i$ centred at the $p_i$,
and that a closed $1$-form $\eta$ on $M$ is exact if and only if its integrals over small circles around all the $p_i$ vanish.
\end{exercise}

\begin{solution}
After a rotation we may assume the $p_i$ have distinct first coordinates, $p_1$ having the largest. Let $U$ be $M$ intersected with $\{x < a\}$ and $V$ be $M$ intersected with
$\{x > b\}$, where $b < a$ are chosen so that $p_1$ is the only point with first coordinate $> b$ and none has first coordinate in $[b, a]$. Then $U \cap V$ is a strip, which is
contractible, so Mayer--Vietoris gives $H^1(M) \cong H^1(U) \oplus H^1(V)$. The half-planes are diffeomorphic to $\RR^2$, so $V$ is a plane minus one point, with $H^1 \cong \RR$ spanned by
$\alpha_1$ (Example~\ref{de-rham:example-angle-form}), and $U$ is a plane minus $k - 1$ points; induction gives $\RR^k$. The integral of $\alpha_i$ over a small circle around $p_j$ is
$2\pi\delta_{ij}$ (away from $p_i$, $\alpha_i$ is exact on a disc), so the map $[\eta] \mapsto (\int_{C_j}\eta)_j$ is surjective onto $\RR^k$, hence an isomorphism, and the $[\alpha_i]$ form a basis.
\end{solution}

\begin{exercise}\label{de-rham:exercise-solid-angle}
On $\RR^3 \setminus \{0\}$ let $\sigma = (x\,dy \wedge dz + y\,dz \wedge dx + z\,dx \wedge dy)/r^3$, $r = (x^2 + y^2 + z^2)^{1/2}$ (the \emph{solid angle form}). Show that $\sigma$ is closed, that
$\int_{S^2}\sigma = 4\pi$, and deduce that $H^2_{\mathrm{dR}}(\RR^3 \setminus \{0\}) = \RR[\sigma]$ and that $\sigma$ is not exact.
\end{exercise}

\begin{solution}
By Example~\ref{differential-forms:example-forms-r3}, $d\sigma = \operatorname{div}(x/r^3, y/r^3, z/r^3)\,dx \wedge dy \wedge dz$, and $\partial_x(x/r^3) = 1/r^3 - 3x^2/r^5$, so the divergence is
$3/r^3 - 3r^2/r^5 = 0$. On the unit sphere $r = 1$ and $\sigma$ restricts to the area form (Example~\ref{differential-forms:example-integrals}(2)), so $\int_{S^2}\sigma = 4\pi$. An exact form would have
integral $0$ over the closed surface $S^2$ (Theorem~\ref{differential-forms:theorem-stokes}). Finally $\RR^3 \setminus \{0\}$ is homotopy equivalent to $S^2$, so its $H^2$ is one-dimensional
(Example~\ref{de-rham:example-spheres}), spanned by the nonzero class $[\sigma]$.
\end{solution}

\begin{exercise}\label{de-rham:exercise-explicit-primitive}
Decide whether the $2$-form $\omega = z\,dx \wedge dy + x\,dy \wedge dz + y\,dz \wedge dx$ on $\RR^3$ is closed, and find a primitive of the closed form
$\omega^{\prime} = dx \wedge dy + dy \wedge dz$ using Corollary~\ref{de-rham:corollary-explicit-poincare}.
\end{exercise}

\begin{solution}
$d\omega = \partial_zz\,dz \wedge dx \wedge dy + \partial_xx\,dx \wedge dy \wedge dz + \partial_yy\,dy \wedge dz \wedge dx = 3\,dx \wedge dy \wedge dz \neq 0$, so $\omega$ is not closed and has no primitive.
For $\omega^{\prime}$, which has constant coefficients and is closed, $\iota_X\omega^{\prime} = x\,dy - y\,dx + y\,dz - z\,dy$ and $\eta = \frac12(x\,dy - y\,dx + y\,dz - z\,dy)$; then
$d\eta = \frac12(2\,dx \wedge dy + 2\,dy \wedge dz) = \omega^{\prime}$.
\end{solution}

\begin{exercise}\label{de-rham:exercise-spheres-euler}
Write $S^n = U \cup V$ with $U$, $V$ the complements of the north and south poles. Using only Proposition~\ref{de-rham:proposition-dimension-counting}, find $\chi(S^n)$. Deduce $b_n(S^n)$ from $b_0 = 1$ and the vanishing of
$b_k(S^n)$ for $0 < k < n$ (Example~\ref{de-rham:example-spheres}), and check the answer against that example.
\end{exercise}

\begin{solution}
$U$ and $V$ are diffeomorphic to $\RR^n$ by stereographic projection, so $\chi(U) = \chi(V) = 1$, and $U \cap V \cong S^{n-1} \times \RR$ is homotopy equivalent to $S^{n-1}$. By induction on $n$, starting from $\chi(S^0) = 2$
(two points), $\chi(S^n) = 2 - \chi(S^{n-1})$, so $\chi(S^n) = 1 + (-1)^n$. Then $1 + (-1)^nb_n = 1 + (-1)^n$ gives $b_n = 1$.
\end{solution}

\begin{exercise}\label{de-rham:exercise-genus-betti}
Let $S$ be a compact connected surface.
\begin{enumerate}
\item Let $D \subseteq S$ be a closed disc in a chart, and $S^\circ = S \setminus D$. Show that $H^\bullet_{\mathrm{dR}}(S^\circ)$ is finite-dimensional and $\chi(S^\circ) = \chi(S) - 1$.
\item The connected sum $S \# S'$ of two compact connected surfaces is obtained by removing a closed disc from each and gluing the two along collars of the boundary circles, so that $S \# S' = U \cup V$ with
$U \cong S^\circ$, $V \cong S'^\circ$ and $U \cap V$ an annulus. Show that $\chi(S \# S') = \chi(S) + \chi(S') - 2$.
\item Deduce that the closed orientable surface $\Sigma_g = T^2 \# \cdots \# T^2$ of genus $g$ (with $\Sigma_0 = S^2$) has $\chi(\Sigma_g) = 2 - 2g$, hence $b_1(\Sigma_g) = 2g$, and that the connected sum $N_h$ of
$h \geq 1$ copies of $\RR\PP^2$ has $\chi(N_h) = 2 - h$, hence $b_1(N_h) = h - 1$. Use that $\Sigma_g$ is orientable and $N_h$ is not, so $b_2(\Sigma_g) = 1$ and $b_2(N_h) = 0$
(Corollary~\ref{de-rham:corollary-consequences}).
\end{enumerate}
\end{exercise}

\begin{solution}
(1) Let $D^\prime \supset D$ be a slightly larger open disc in the same chart. Then $S = S^\circ \cup D^\prime$, and $S^\circ \cap D^\prime = D^\prime \setminus D$ is an open annulus, homotopy equivalent to $S^1$, so
$\chi(S^\circ \cap D^\prime) = 0$. As $S$ and the annulus have finite-dimensional cohomology (Corollary~\ref{de-rham:corollary-consequences}), so does $S^\circ \sqcup D^\prime$ by
Proposition~\ref{de-rham:proposition-dimension-counting}(2), and $\chi(S) = \chi(S^\circ) + \chi(D^\prime) - 0 = \chi(S^\circ) + 1$.

(2) By (1) and Proposition~\ref{de-rham:proposition-dimension-counting}(2), $\chi(S \# S') = (\chi(S) - 1) + (\chi(S') - 1) - 0$.

(3) $\chi(T^2) = 0$ (Example~\ref{de-rham:example-dimension-counting}) and $\chi(S^2) = 2$, so by (2) $\chi(\Sigma_g) = \chi(\Sigma_{g-1}) - 2$ and $\chi(\Sigma_g) = 2 - 2g$; then
$b_1 = b_0 + b_2 - \chi = 2 - (2 - 2g) = 2g$. For $\RR\PP^2$, Example~\ref{de-rham:example-real-projective} gives $\chi = 1$, so $\chi(N_h) = \chi(N_{h-1}) - 1 = 2 - h$ and $b_1(N_h) = 1 + 0 - (2 - h) = h - 1$. For instance the Klein bottle
$N_2$ has $b_1 = 1$, as in Example~\ref{de-rham:example-dimension-counting}(2). The integral homology is $H_1(\Sigma_g) \cong \ZZ^{2g}$ and $H_1(N_h) \cong \ZZ^{h-1} \oplus \ZZ/2\ZZ$ \cite{hatcher-at}: the torsion of the
non-orientable surfaces is again invisible.
\end{solution}

\begin{exercise}\label{de-rham:exercise-dimension-in-sequences}
Let $0 \to A \to B \to C \to 0$ be a short exact sequence of finite-dimensional real vector spaces. Show that $B \cong A \oplus C$, so that over a field every short exact sequence splits. Give an example of a short exact
sequence of abelian groups that does not split, and explain why such sequences cannot occur among the de Rham cohomology groups of a Mayer--Vietoris sequence, but do occur in singular homology with
coefficients in $\ZZ$.
\end{exercise}

\begin{solution}
Let $a \colon A \to B$ and $b \colon B \to C$ be the maps. Choose a basis $c_1, \ldots, c_r$ of $C$ and, as $b$ is surjective, elements $b_i$ with $b(b_i) = c_i$; the linear map $s \colon C \to B$ with $s(c_i) = b_i$ satisfies
$b \circ s = \id_C$. Then $A \oplus C \to B$, $(x, y) \mapsto a(x) + s(y)$, is an isomorphism: it is injective, since $a(x) + s(y) = 0$ gives $y = b(a(x) + s(y)) = 0$ and then $x = 0$ as $a$ is injective; and it is surjective, since
$z - s(b(z)) \in \ker b = \im a$ for every $z \in B$. This argument works for vector spaces over any field, and for modules wherever $C$ has a basis.

The sequence $0 \to \ZZ \xrightarrow{2} \ZZ \to \ZZ/2\ZZ \to 0$ does not split: $\ZZ \not\cong \ZZ \oplus \ZZ/2\ZZ$, since $\ZZ$ has no element of order $2$. In de Rham cohomology all groups are real vector spaces and all
maps are linear, so every short exact sequence extracted from a long exact sequence splits, and only dimensions matter. In integral homology such non-split sequences appear: in the cellular chain complex of
$\RR\PP^2$ the boundary of the $2$-cell is twice the $1$-cell, and $H_1(\RR\PP^2) \cong \ZZ/2\ZZ$ arises exactly from the multiplication by $2$ above.
\end{solution}
